%=====================================================================
%  FICHE DE RÉVISION — Réseaux de neurones pour les EDP (PINN & Deep Ritz)
%  Projet 1A — CERMICS.  Document autonome.
%  Compilation : pdflatex (2 passes pour le sommaire)
%=====================================================================
\documentclass[11pt]{article}
\usepackage[utf8]{inputenc}
\usepackage[T1]{fontenc}
\IfFileExists{french.ldf}{\usepackage[french]{babel}}{}
\usepackage{lmodern}
\usepackage{amsmath,amssymb}
\usepackage{xcolor}
\usepackage[a4paper,margin=2.1cm]{geometry}
\usepackage{titlesec}
\usepackage{enumitem}
\usepackage{booktabs}
\usepackage{listings}
\usepackage[most]{tcolorbox}
\usepackage{microtype}
\usepackage[hidelinks]{hyperref}

\definecolor{bblue}{RGB}{51,51,178}
\definecolor{dblue}{RGB}{38,38,134}
\definecolor{lblue}{RGB}{233,233,243}
\definecolor{codebg}{RGB}{237,242,234}

\titleformat{\section}{\sffamily\large\bfseries\color{bblue}}{\thesection.}{0.5em}{}
\titleformat{\subsection}{\sffamily\bfseries\color{dblue}}{\thesubsection}{0.5em}{}
\titlespacing*{\section}{0pt}{10pt}{4pt}
\setlist[itemize]{leftmargin=1.3em,itemsep=1pt,topsep=2pt}

% Boîte "formule clé"
\newtcolorbox{formbox}[1][]{colback=lblue,colframe=bblue,boxrule=0.7pt,arc=2pt,
  left=8pt,right=8pt,top=4pt,bottom=4pt,#1}
% Boîte "question de jury"
\newtcolorbox{qbox}[1]{colback=white,colframe=dblue,boxrule=0.7pt,arc=2pt,
  left=8pt,right=8pt,top=4pt,bottom=4pt,
  title={\sffamily\bfseries #1},coltitle=white,colbacktitle=dblue,fonttitle=\bfseries}
\newcommand{\key}[1]{\textbf{#1}}

\lstset{language=Python,basicstyle=\ttfamily\footnotesize,backgroundcolor=\color{codebg},
  keywordstyle=\color{bblue}\bfseries,commentstyle=\color{gray}\itshape,
  stringstyle=\color{dblue},showstringspaces=false,breaklines=true,
  frame=single,rulecolor=\color{bblue!40},framesep=4pt,xleftmargin=4pt,
  literate={é}{{\'e}}1 {è}{{\`e}}1 {à}{{\`a}}1 {ô}{{\^o}}1 {É}{{\'E}}1 {â}{{\^a}}1 {ç}{{\c c}}1 {î}{{\^i}}1 {û}{{\^u}}1}

\setlength{\parindent}{0pt}
\setlength{\parskip}{3pt}

\begin{document}

\begin{center}
  {\sffamily\bfseries\color{dblue}\Large Fiche de révision — Maîtrise du sujet}\\[3pt]
  {\sffamily\color{bblue}\large Réseaux de neurones pour la résolution d'EDP : PINN \& Deep Ritz}\\[5pt]
  {\small Projet 1A \;\textbullet\; CERMICS \;\textbullet\; Paul Rivaud \;\textbullet\; juin 2026}
\end{center}
\hrule height 0.8pt
\vspace{4pt}

{\small\tableofcontents}
\vspace{4pt}\hrule height 0.4pt

%==================================================================
\section{Vue d'ensemble}
Le projet résout numériquement deux EDP 1D sur l'ouvert $\Omega=\,]0,1[$ —
l'\key{équation de Poisson} et l'\key{équation de Helmholtz modifiée} — en
\key{approchant la solution par un réseau de neurones}. Deux paradigmes sont
implémentés en PyTorch :
\begin{itemize}
  \item \key{PINN} (\emph{Physics-Informed Neural Network}) : formulation
        \key{forte}, on minimise le \key{résidu} de l'EDP ;
  \item \key{Deep Ritz} : formulation \key{faible/variationnelle}, on minimise une
        \key{énergie}.
\end{itemize}
Idée commune : transformer « résoudre une EDP » en « \key{minimiser une fonction de
perte} » par descente de gradient, les dérivées du réseau étant obtenues par
\key{différentiation automatique}.

%==================================================================
\section{Concepts clés}

\subsection*{EDP : formulation forte vs faible}
\begin{itemize}
  \item \key{Forte} : on cherche $u\in C^2(\Omega)$ vérifiant l'équation
        \emph{en tout point}. Nécessite ici une \key{dérivée seconde}.
  \item \key{Faible (variationnelle)} : on cherche $u\in H^1_0(\Omega)$ tel que
        $a(u,v)=b(v)$ pour toute fonction test $v$. Ne demande qu'une
        \key{dérivée première} ; cadre du théorème de \key{Lax-Milgram}
        (existence \& unicité si $a$ continue et coercive).
\end{itemize}

\subsection*{Réseau de neurones}
Modèle paramétrique : un graphe de neurones organisés en couches, chacun calculant
une \key{transformation non linéaire} d'une combinaison linéaire pondérée de ses
entrées. La non-linéarité vient de la \key{fonction d'activation} $\sigma$.
\key{Apprendre} = minimiser une perte par descente de gradient + \key{rétropropagation}.

\subsection*{PINN}
Le réseau $u_\theta$ \key{est} la solution approchée. On injecte l'EDP et les
conditions aux bords (CL) dans la perte via le \key{résidu}. Cadre
\key{non supervisé} : pas de données $(x,u(x))$, la « supervision » vient de la
physique. Convient aux problèmes \key{directs et inverses}.

\subsection*{Fonction de perte physique}
Somme d'un terme \key{intérieur} (résidu de l'EDP sur des \key{points de
collocation}) et d'un terme de \key{bord} (pénalisation des CL).

\subsection*{Entraînement}
\key{Différentiation automatique} (dérivées exactes du réseau, sans pas $h$ ni
erreur de troncature) ; \key{descente de gradient} (optimiseur \key{Adam}) ;
intégrales approchées par la \key{méthode des rectangles}.

%==================================================================
\section{Formules essentielles}

\textbf{Problèmes test.}
\[
  \text{(Poisson)}\quad
  \begin{cases}-u''(x)=f(x), & x\in\,]0,1[\\ u(0)=u_0,\ u(1)=u_1\end{cases}
  \qquad\qquad
  \text{(Helmholtz)}\quad -\Delta u+\lambda u=f,\ \ \lambda>0.
\]

\textbf{Formulation faible (Poisson)} : trouver $u\in H^1_0(0,1)$ tel que
$a(u,v)=b(v)\ \forall v\in H^1_0$, avec
\[
  a(u,v)=\int_0^1 u'(x)\,v'(x)\,dx,\qquad b(v)=\int_0^1 f(x)\,v(x)\,dx .
\]

\textbf{Espace des réseaux \& approximation universelle (Cybenko, 1989).}
\begin{formbox}
\[
  V_{nn}=\Bigl\{\,x\mapsto \textstyle\sum_{i=1}^{n} c_i\,\sigma(w_i x+b_i)\ \Big|\
  \theta=(c_i,w_i,b_i)_{1\le i\le n}\in\mathbb{R}^{3n}\Bigr\}.
\]
Un réseau à \emph{une} couche cachée + activation continue approxime
\emph{uniformément} toute fonction continue.
\end{formbox}

\textbf{Méthode PINN — fonctionnelle d'énergie (formulation forte).}
\begin{formbox}
\[
  E(u_\theta)=\int_{\Omega}\bigl|u_\theta''(x)+f(x)\bigr|^2dx
   +\bigl(u_\theta(0)-u_0\bigr)^2+\bigl(u_\theta(1)-u_1\bigr)^2,
  \qquad E(u^\star)=0 .
\]
\end{formbox}
Discrétisation par collocation sur $N$ points $x_i$ (perte effectivement minimisée) :
\[
  \mathcal{L}(\theta)=
  \underbrace{\frac1N\sum_{i=1}^{N}\bigl(u_\theta''(x_i)+f(x_i)\bigr)^2}_{\mathcal{L}_{\mathrm{PDE}}}
  +\underbrace{\bigl(u_\theta(0)-u_0\bigr)^2+\bigl(u_\theta(1)-u_1\bigr)^2}_{\mathcal{L}_{\mathrm{BC}}},
  \qquad \theta^\star=\arg\min_\theta \mathcal{L}(\theta).
\]

\textbf{Méthode Deep Ritz — énergie (formulation faible).} Si $a$ symétrique,
la solution minimise $J(v)=\tfrac12 a(v,v)-b(v)$. Cas concrets (avec pénalisation
des bords $\alpha,\beta$) :
\begin{formbox}
\[
  J_{\text{Poisson}}(u)=\frac12\!\int_\Omega |u'|^2dx-\!\int_\Omega f\,u\,dx,
  \qquad
  J_{\text{Helmholtz}}(u)=\frac12\!\int_\Omega |u'|^2dx+\frac{\lambda}{2}\!\int_\Omega |u|^2dx-\!\int_\Omega f\,u\,dx .
\]
\end{formbox}

\textbf{Activation \& optimisation.}
\[
  \sigma(x)=\tanh(x),\qquad \sigma'(x)=1-\tanh^2(x),\qquad
  \theta \leftarrow \theta-\eta\,\nabla_\theta\mathcal{L}(\theta)\ \ (\text{Adam}).
\]
$\tanh\in C^\infty\Rightarrow u_\theta\in C^2$ et $u_\theta''$ calculable par autodiff.

%==================================================================
\section{L'implémentation (le code)}
\textbf{Pile technique :} PyTorch + \texttt{torch.func} (\texttt{grad}, \texttt{vmap},
\texttt{functional\_call}).

\textbf{Architectures.} \texttt{snn} : réseau à \key{une} couche cachée
($3n$ paramètres) pour PINN. \texttt{NN} : réseau \key{profond}
(\texttt{ModuleList}, profondeur variable) pour Deep Ritz.

\textbf{Dérivées par autodiff} (point central) : on dérive \emph{par rapport à
l'entrée} $x$, deux fois, par composition de \texttt{grad}.
\begin{lstlisting}
def f_forward(x, net, params, buffers):
    return torch.func.functional_call(net, (params, buffers), x.unsqueeze(-1)).squeeze()

f_forward_x  = grad(f_forward, argnums=0)     # u'_theta  (derivee 1ere)
f_forward_xx = grad(f_forward_x, argnums=0)   # u''_theta (derivee 2nde)
\end{lstlisting}

\textbf{Boucle PINN} (\texttt{optimal\_snn}) : $N=1000$ points, Adam $\eta=0{,}01$,
1000 époques ; perte = résidu PDE + carrés des valeurs aux bords ; vectorisation
par \texttt{vmap}.
\begin{lstlisting}
loss  = ((f_forward_xx_grid - f_grid)**2).mean()      # L_PDE  (residu)
loss += f_forward(bc_0, net, params, buffers)**2       # L_BC en x=0
loss += f_forward(bc_1, net, params, buffers)**2       # L_BC en x=1
optimizer.zero_grad(); loss.backward(); optimizer.step()
\end{lstlisting}

\textbf{Version optimisée} (\texttt{snn\_optimized} / \texttt{optimal\_snn\_optimized}) :
choix de l'activation via un dictionnaire (\texttt{tanh}, \texttt{relu},
\texttt{sigmoid}), \key{jeu de validation} séparé, \key{sauvegarde des meilleurs
poids} (\texttt{val\_loss} minimale, \texttt{.clone()}) pour contrer le
sur-apprentissage, et calcul de l'\key{erreur $L^2$} contre la solution exacte.

\textbf{Deep Ritz} (\texttt{DeepRitz}, \texttt{DeepRitz2}) : on n'utilise que la
\key{dérivée première}, plus une \key{pénalisation des bords} $\alpha,\beta$.
\begin{lstlisting}
# Poisson : J = 0.5 <|u'|^2> - <f u> + penalites de bord
loss = 0.5*(u_x_grid**2).mean() - (u_grid*f_grid).mean() \
       + 10*f_forward(bc_0,net,params,buffers)**2 \
       + 10*f_forward(bc_1,net,params,buffers)**2
\end{lstlisting}

\textbf{Choix techniques en bref :} dérivées \key{exactes} par autodiff (vs
différences finies) ; \texttt{tanh} pour la régularité $C^\infty$ ; \texttt{vmap}
pour vectoriser sans boucle \texttt{for} ; \texttt{Adam} (pas adaptatif) ;
validation + meilleur modèle sauvegardé.

%==================================================================
\section{Résultats clés à retenir}
\begin{center}\small
\begin{tabular}{llccr}
\toprule
Configuration & Activation & Neurones & Couches & Loss val.\ min.\\
\midrule
PINN de base   & tanh    & 3  & 1  & $5{,}13\times10^{-3}$\\
PINN optimisé  & tanh    & 5  & 1  & $1{,}03\times10^{-4}$\\
PINN optimisé  & tanh    & 10 & 1  & $1{,}71\times10^{-5}$\\
PINN optimisé  & tanh    & 50 & 1  & $1{,}80\times10^{-5}$\\
PINN optimisé  & sigmoid & 50 & 1  & $1{,}58\times10^{-5}$\\
PINN optimisé  & ReLU    & 50 & 1  & $4{,}99\times10^{-1}$\\
DNN stable     & tanh    & 5  & 13 & $\approx 10^{-2}$\\
DNN instable   & tanh    & 5  & 20 & $\approx 4{,}8\times10^{-1}$\\
\bottomrule
\end{tabular}
\end{center}
\begin{itemize}
  \item \key{tanh} est le meilleur choix ; \key{10 neurones suffisent} en 1D
        (gain négligeable au-delà).
  \item \key{ReLU échoue} : sa dérivée seconde est nulle p.p., le résidu
        $u_\theta''+f$ ne peut pas être minimisé.
  \item \key{Profondeur} : au-delà de $\sim$17 couches, la perte explose
        (\key{vanishing gradient} : facteurs $1-\tanh^2\!\approx 0$). Le problème
        1D ne justifie pas un réseau profond.
  \item \key{Deep Ritz} converge plus lentement que PINN sur Poisson, mais évite
        la dérivée seconde et s'applique quand l'énergie PINN n'est pas évidente
        (cas Helmholtz, convergence rapide, $L^2\!\sim\!10^{-5}$).
  \item \key{Pénalisation $\alpha,\beta$} (Deep Ritz) : trop faible $\Rightarrow$
        \emph{underfitting} des bords ; trop forte $\Rightarrow$ les CL écrasent
        l'EDP intérieure (\emph{overfitting} des bords).
\end{itemize}

%==================================================================
\section{Applications concrètes des PINN}
Au-delà du cas 1D du projet, les PINN sont utilisés pour :
\begin{itemize}
  \item \key{mécanique des fluides} (Navier-Stokes), aérodynamique, transferts
        thermiques, diffusion ;
  \item \key{ondes \& électromagnétisme} (Helmholtz, Maxwell), élasticité ;
  \item \key{problèmes inverses} : identifier des paramètres / sources / propriétés
        matériau à partir de mesures partielles ;
  \item \key{assimilation de données} : combiner données rares et lois physiques ;
  \item \key{grande dimension} : EDP où le maillage est impraticable
        (finance — Black-Scholes, équations de Fokker-Planck\ldots).
\end{itemize}
Atouts : \key{sans maillage}, gèrent des géométries complexes, unifient
problèmes directs et inverses.

%==================================================================
\section{Limites \& ouvertures}
\subsection*{Limites}
\begin{itemize}
  \item \key{Optimisation non convexe} : pas de garantie d'atteindre le minimum
        global (Cybenko garantit l'\emph{existence} d'un bon approximant, pas que la
        descente le trouve).
  \item \key{Équilibrage de la perte} : le poids relatif PDE / bords (et $\lambda$,
        $\alpha$, $\beta$) est délicat à régler.
  \item \key{Coût} des dérivées d'ordre élevé (double autodiff pour $u''$) ;
        convergence parfois lente.
  \item En \key{petite dimension}, souvent moins précis/rapide que les éléments finis.
  \item \key{Capacité} limitée $\Rightarrow$ plateau de la loss ; profondeur
        $\Rightarrow$ \key{vanishing gradient} avec tanh.
\end{itemize}
\subsection*{Ouvertures}
\begin{itemize}
  \item \key{Deep Ritz} : formulation faible, ne demande que $u'$, applicable plus
        largement.
  \item \key{Conditions aux bords « dures »} (imposées par construction) plutôt que
        par pénalisation.
  \item \key{Pondération adaptative} des termes de la perte ; échantillonnage
        adaptatif des points de collocation.
  \item \key{Architectures dédiées} (\emph{Fourier features}, SIREN) pour les
        hautes fréquences ; extension \key{2D/3D}, équations non linéaires et
        instationnaires, problèmes inverses.
\end{itemize}

%==================================================================
\section{Questions de jury probables \& réponses}

\begin{qbox}{Pourquoi \texttt{tanh} et pas \texttt{ReLU} ?}
ReLU a une dérivée seconde \key{nulle presque partout} : le terme $u_\theta''$ du
résidu est alors structurellement nul, la perte PDE ne peut pas être minimisée.
\texttt{tanh} est $C^\infty$, donc $u_\theta''$ est bien définie et non nulle.
\end{qbox}

\begin{qbox}{Différence entre PINN et Deep Ritz ?}
PINN : formulation \key{forte}, on minimise le résidu (il faut $u''$). Deep Ritz :
formulation \key{faible}, on minimise une énergie (il ne faut que $u'$) — plus
économique, et applicable même quand écrire une énergie « PINN » n'est pas évident
(cas Helmholtz).
\end{qbox}

\begin{qbox}{Qu'est-ce que la différentiation automatique ? Différence avec les différences finies ?}
C'est le calcul \key{exact} des dérivées via la règle de la chaîne appliquée au
graphe de calcul du réseau. Contrairement aux différences finies, \key{pas de pas
$h$ à choisir ni d'erreur de troncature}.
\end{qbox}

\begin{qbox}{Pourquoi parle-t-on d'apprentissage « non supervisé » ?}
Il n'y a \key{aucune donnée} $(x,u(x))$. La « cible » est l'équation elle-même :
on minimise son résidu. La physique remplace les étiquettes.
\end{qbox}

\begin{qbox}{Comment les conditions aux bords sont-elles imposées ?}
Par \key{pénalisation} dans la perte (contrainte « molle »). Alternative possible :
les imposer \key{par construction} (contrainte « dure »), p.\,ex.\ via un ansatz
qui annule l'erreur de bord.
\end{qbox}

\begin{qbox}{Rôle des coefficients $\alpha,\beta$ (Deep Ritz) ?}
Ils pondèrent la pénalisation des bords. Trop faibles : \key{underfitting} des CL.
Trop forts : l'algorithme « oublie » l'EDP à l'intérieur (\key{overfitting} des bords).
\end{qbox}

\begin{qbox}{Pourquoi la loss stagne-t-elle (plateau) avec 3 neurones ?}
\key{Capacité} insuffisante de $V_{nn}$ : le résidu ne peut pas être annulé.
Il faut davantage de neurones (10 suffisent ici).
\end{qbox}

\begin{qbox}{Pourquoi un réseau très profond se dégrade-t-il ?}
\key{Vanishing gradient} : en profondeur, les gradients sont multipliés par
$1-\tanh^2\!\approx 0$, la rétropropagation devient inefficace. De plus, le
problème 1D n'a pas de structure hiérarchique justifiant la profondeur.
\end{qbox}

\begin{qbox}{A-t-on une garantie de convergence vers la vraie solution ?}
Non : l'optimisation est \key{non convexe}. On a une garantie d'\key{existence}
d'un bon approximant (Cybenko), mais pas que la descente de gradient atteigne le
minimum global.
\end{qbox}

\begin{qbox}{Combien de paramètres a le réseau à une couche ?}
$3n$ : pour chaque neurone $i$, le poids $w_i$, le biais $b_i$ et le coefficient de
sortie $c_i$.
\end{qbox}

\begin{qbox}{À quoi sert \texttt{vmap} ?}
À \key{vectoriser} l'évaluation du réseau et de ses dérivées sur tous les points de
collocation, sans boucle Python — donc beaucoup plus rapide.
\end{qbox}

\begin{qbox}{Que mesure $E(u)$, et pourquoi $E(u^\star)=0$ ?}
Le \key{résidu intérieur} (violation de l'EDP) + la \key{violation des CL}. La
solution exacte annule les deux termes, d'où $E(u^\star)=0$.
\end{qbox}

\begin{qbox}{À quoi sert le théorème de Lax-Milgram ?}
Il garantit l'\key{existence et l'unicité} de la solution faible quand la forme
bilinéaire $a$ est continue et \key{coercive} — fondement du cadre variationnel
(Deep Ritz).
\end{qbox}

\begin{qbox}{La solution exacte $\sin(2\pi x)$ sert-elle à l'entraînement ?}
Non : elle ne sert qu'à \key{valider} (calcul de l'erreur $L^2$). L'entraînement
n'utilise que $f$ et les conditions aux bords.
\end{qbox}

\vfill
\hrule height 0.4pt
{\footnotesize\color{gray}\textbf{Références :} Cybenko (1989), \emph{Approximation by
superpositions of a sigmoidal function} ; Weinan E \& Yu (2018), \emph{The Deep Ritz
Method} ; Raissi, Perdikaris \& Karniadakis (2019), \emph{Physics-Informed Neural
Networks}, J.\ Comput.\ Phys.\ 378.}

\end{document}
